\begin{figure}[htbp]
\centering
\begin{tikzpicture}
\begin{axis}[
  width=0.78\textwidth, height=6.2cm,
  xlabel={Irradiation time [days]},
  ylabel={Gap width [\si{\micro\metre}]},
  xmin=0, xmax=760, ymin=-11, ymax=19,
  grid=major, grid style={cGrid!60,very thin},
  legend pos=north east, legend style={font=\scriptsize,draw=black!30},
  label style={font=\small}, tick label style={font=\scriptsize},
]
\addplot[color=cClad,thick,no marks] coordinates {
 (69.4,16.312) (100.9,14.651) (132.3,12.955) (163.7,11.242) (195.1,9.527)
 (226.5,7.828) (257.9,6.162) (289.4,4.544) (320.8,2.991) (352.2,1.517)
 (383.6,0.134) (415.0,-1.147) (446.4,-2.319) (477.8,-3.376) (509.3,-4.313)
 (540.7,-5.129) (572.1,-5.825) (603.5,-6.401) (634.9,-6.863) (666.3,-7.215)
 (697.8,-7.466) (729.2,-7.621)};
\addlegendentry{Surrogate (\SI{2}{\milli\second})}
\addplot[only marks,mark=*,mark size=2.2pt,color=cFuel] coordinates {
 (170.0,10.097) (354.1,0.370) (456.9,-1.959)
 (590.2,-5.251) (710.9,-7.703) (729.2,-7.997)};
\addlegendentry{OFFBEAT (\SI{36}{\second})}
\addplot[black!55,dashed,thick,domain=0:760,forget plot] {0};
\node[font=\tiny,anchor=south west,text=black!65] at (axis cs:20,0.6)
  {gap open};
\node[font=\tiny,anchor=north west,text=black!65] at (axis cs:20,-0.6)
  {contact established};
\end{axis}
\end{tikzpicture}
\caption[PCMI dating at an unseen power]{Pellet-cladding gap closure at
\SI{25}{\kilo\watt\per\metre}, a power absent from the surrogate's training set.
The continuous curve is the surrogate prediction, the markers are the states
computed by OFFBEAT. The zero crossing, the onset of PCMI, is predicted at
\num{387} days against \num{370} computed, a deviation of
\SI{4.4}{\percent}.}
\label{fig:pcmi}
\end{figure}

The surrogate places gap closure at \textbf{\num{387} days}, against
\textbf{\num{370} days} obtained by interpolating the states computed by
OFFBEAT: a deviation of \SI{4.4}{\percent}, that is sixteen days over a two-year
cycle. The cost of the prediction is \SI{2}{\milli\second}, against
\SI{36}{\second} for the complete simulation, an acceleration factor of about
\num{18000}.

This result validates the surrogate where it matters: no longer on its ability
to reproduce an equilibrium temperature, but on its ability to \emph{date a
safety event} by interpolating between operating conditions it has never seen.

Two reservations nonetheless remain. The validation rests on twenty-five points,
which is still few; and nothing guarantees validity \textbf{outside} the training domain; a surrogate extrapolates poorly by construction, as the first
version demonstrated spectacularly. This is precisely why the uncertainty
provided by the Gaussian process is systematically reported at the interface
alongside the predicted value.

\section{Stress-testing self-healing}
\label{sec:epreuve}

The aim of this section is to delimit, in a measured rather than anecdotal way,
what the agent can repair on its own.

\subsection{Protocol}

I built a benchmark on the principle of \textbf{fault injection}. The creation
tool generates a healthy case, I corrupt it in one controlled way, and the
executor then receives it with self-healing enabled.
Fourteen cases were processed: twelve faults spread over four root-cause
categories, plus two healthy cases serving as controls.

The outcome of each case is classified along two deliberately
\textbf{independent} axes:
\begin{itemize}
  \item \textbf{detected}, did a pattern from the knowledge base recognise the
    error, as opposed to falling back on the language model;
  \item \textbf{repaired}, did the case eventually run.
\end{itemize}

Separating these two axes is the point of the protocol: it is the ``detected but
not repaired'' cell that materialises the gap between a \emph{symptom}
recognised in a log and a \emph{cause} actually addressed.

Each fault is further labelled a priori according to two distinct criteria: does
the base offer a correction for this pattern, and can that correction, by its
nature, address the root cause? A fault for which the base \emph{believes} it
can intervene while the remedy targets only the symptom is the heart of the
problem studied here.

\subsection{Results}

Table~\ref{tab:selfheal} gives the results by category.

\begin{table}[htbp]
\centering
\caption[Self-healing benchmark]{Self-healing benchmark: 14 cases, of which 12
injected faults. The ``stopped at meshing'' column marks the cases intercepted
by the mesher before the solver was even launched.}
\label{tab:selfheal}
\begin{tabularx}{\textwidth}{@{}Xcccc@{}}
\toprule
\textbf{Root-cause category} & \textbf{$n$} & \textbf{detected} &
\textbf{repaired} & \textbf{stopped at meshing}\\
\midrule
Control (healthy cases) & 2 & 0/2 & \textbf{2/2} & 0/2\\
Numerical (solver divergence) & 3 & 3/3 & 0/3 & 0/3\\
Physical consistency (geometry) & 3 & 3/3 & 0/3 & 2/3\\
Missing file & 3 & 1/3 & 0/3 & 0/3\\
Invalid dictionary entry & 3 & 2/3 & \textbf{1/3} & 1/3\\
\midrule
\textbf{Total} & \textbf{14} & \textbf{9} & \textbf{3} & \textbf{3}\\
\bottomrule
\end{tabularx}
\end{table}

The detection rate is reasonable: nine faults out of twelve are recognised by a
pattern, the language model fallback having been called upon only twice. The
repair rate, on the other hand, is very low: of the twelve injected faults,
\textbf{only one} was repaired.

The cross-tabulation of table~\ref{tab:selfheal-croise} explains why, and
constitutes the central result of this section.

\begin{table}[htbp]
\centering
\caption[Repair versus treatability of the cause]{Repair rate depending on
whether the available correction can, or cannot, address the root cause of the
fault.}
\label{tab:selfheal-croise}
\begin{tabularx}{0.9\textwidth}{@{}Xc@{}}
\toprule
\textbf{Situation} & \textbf{Repaired}\\
\midrule
Correction in the base \textbf{and} cause treatable by it & \textbf{1/1}\\
Correction in the base but cause \textbf{not} treatable by it &
  \textcolor{cAlert}{\textbf{0/3}}\\
No correction in the base for this pattern & 0/8\\
\bottomrule
\end{tabularx}
\end{table}

Self-healing therefore repairs \emph{exactly} the faults whose root cause falls
within the scope of the available correction, and only those. The middle row is
the most instructive: in those three cases the pattern is recognised, the
diagnosis is displayed, the correction fires, and fails every time, consuming
the full budget of three attempts each time. The mechanism gives the appearance
of working while it structurally cannot succeed.

\subsection{Analysis of the failures}

Three distinct failure mechanisms emerge.

\textbf{The remedy targets only the symptom.} The three numerical faults (powers
\num{100} to \num{200} times nominal, or negative) trigger a floating-point
exception in the creep model. The pattern matches, and the associated correction
reduces the time step. But no reduction of the time step can fix a physically
absurd input: the cause lies upstream of the solver. The agent therefore applies
an ineffective remedy three times in a row.

\textbf{The correction is undone by the solver.} One case is particularly
revealing: the time-step reduction is indeed applied and the first computation
step succeeds, but OFFBEAT's adaptive time-step controller, seeing that all is
going well, immediately increases the step again and causes a new divergence. A
robust repair should act on the \emph{maximum} allowed time step, not on the
initial one.

\textbf{The correction converges too slowly for its budget.} The benchmark
brought to light a defect that one-off trials had not revealed. The correction
associated with exceeding the power history reduced the simulated duration by
\SI{1}{\percent} per attempt. To bring a duration of \SI{1.5e8}{\second} below a
table ending at \SI{1.26e8}{\second} would have required \textbf{eighteen}
attempts for a budget of \textbf{three}: the correction could mathematically
never succeed. It was replaced by a deterministic computation, which reads the
time tables actually present in the case and targets an admissible value
directly. After correction, this case, the only one whose root cause did fall
within an available correction, goes from failure to success in two attempts.

\subsection{A nuance about the ``successful'' failures}

The ``stopped at meshing'' column of table~\ref{tab:selfheal} deserves to be
read against the grain. Two of the three geometric faults are intercepted by the
mesher \emph{before} the solver starts: \texttt{blockMesh} refuses to build an
impossible geometry. This is not a failure of self-healing, it is a safeguard
doing its job. Counting these cases as failures, as a first version of the
benchmark did, amounted to penalising the system for correctly rejecting an
invalid input; hence the addition of this column.

\subsection{Lesson}

Pattern-based self-healing is suited to errors that are \emph{genuinely}
numerical, for which adjusting a resolution parameter makes sense. It is
structurally unsuited to \textbf{physical consistency} errors and to
\textbf{configuration} errors, which account for eleven of the twelve faults in
this benchmark.

The practical consequence is a shift of effort: rather than endlessly enriching
the base of crash patterns, it is more profitable to \textbf{validate the inputs
before execution}. A case that cannot be physically valid should never reach the
solver. A first geometric validation was added to that end (§\ref{sec:bugs}).
